\documentclass[11pt]{article}
\usepackage{amsmath}
\title{When redundant receipts crowd out a different duty}
\date{27 November 2026}
\begin{document}
\maketitle
\paragraph{Audit.}
One ordered native SubT development replay executes two exact
old scoring results and their returning receipts. Distinct native
artifact results \(A,B\) enter a producer-local SQLite archive
only when the exact key matches the required report. Real broker
callbacks stage \(A\) at each named read-only relay \(P,Q\), and
\(B\) at the different relay \(R\). The same scorer process then
executes independent work that closes the lookup. A compatible
second fault can remove either \(A\) relay and its stored copy.
The old score effects are already settled; confirmation of each
of the two results is a different still-open duty. Selection,
closure timing and authored source scope are fixed. This is not
a field or independent mission result.

\paragraph{A larger cap can lower service under a fixed order.}
The selected native application records each named delivered
result at most once. That does not refund receive bytes for
a duplicate that made it through transport. With a three-byte
\(A\), a two-byte \(B\), both \(A\) relays alive, and order
\(A,A,B\), the five-byte limit refuses the second \(A\) at
receive accounting, leaving room for \(B\). At six or seven,
the two copies of \(A\) fit and \(B\) no longer does. At eight,
all fit. Removing one \(A\) relay hides the six-byte failure.
In the five- and six-byte runs all eight intended bytes
were charged on the three sender ledgers. The amount charged
is not a count of fulfilled duties. More capacity has not ruled
out a better policy; it changes the behavior of the fixed one.

For payload lengths \(a,b>0\) and common limit \(C\), assuming
both first packets can transmit, healthy duplicate-first
rejects \(B\) in the band \(2a\le C<2a+b\).
The surprising part is the overlap with \(a+b\le C\), where
one copy of each would fit. If \(a+b\le C<2a\),
the refusing second copy leaves exactly the needed space.
Putting \(B\) before the second \(A\), with the same byte
sizes, makes both old duties timely for any \(C\ge a+b\)
under either single \(A\)-relay loss or no loss in this fixture.
It does not recover a missing \(B\)-relay copy. This ordering
is a scoped admission requirement, not a ban on redundancy.

\paragraph{Current sender limits and available feedback matter.}
The receiver uses the limit of the current sender. With
\(a=2,\ b=3\) and sender limits \(4,4,5\) on \(P,Q,R\),
each possible one-\(A\)-loss case is individually feasible.
But no one fixed order handles both single losses and
the healthy case. Each sole surviving \(A\) needs to go
before \(B\), since its limit is 4 and it cannot follow
\(B\)'s three charged bytes. Putting both before \(B\)
starves \(B\) when both are healthy. The separate executable
alternative sends \(P\) first, skips the backup only after
the current \(A\) callback is actually observed, and otherwise
gives \(Q\) a chance before \(B\). This meets both duties in
the three cases. It assumes the local callback completes at
the dispatch cut before the next opportunity; it does not
apply if receiver bytes can be charged without making that
callback visible. Absence of a delivered notification is not
generally absence of transport effect.

\paragraph{Comparison.}
An independently implemented direct ledger rule with the same
available copies, sender limits, permitted order and callback
cut matches all 4,368 scoped native cases. The third native
score and separate radio work execute throughout. Counts are
grid checks, not a probability estimate. Four Lean arithmetic
lemmas check the admission facts, not the source correspondence.
The executable controls distinguish two individually useful
read handlers from incompatible joint receive use. This
does not establish a gain over a matched direct scheduler,
break a physical precedent, or close I1/I2.
\end{document}
